\section{Experimental Setup}
\label{sec:setup}

\subsection{Design families and frozen split}

The benchmark contains streaming FIR, reverse-FIR, polynomial, recursive-filter,
and median-filter accelerators under a uniform clocked interface.
Specifications and executable reference functions are generated
programmatically.  The resulting domain spans \claim{FamilyCount} families and
several implementation templates, but remains bounded by those generators.

The held-out split was committed before the policy used for the primary result
was trained.  Interpolation contains \claim{InterpDesignCount} unseen parameter
choices inside training ranges.  Extrapolation contains
\claim{ExtrapDesignCount} larger or newly parameterized designs outside those
ranges.  No held-out design appears in supervised examples, policy prompts, or
proxy-training rows.

\subsection{Models, sampling, and functional evaluation}

The primary comparison uses RTLCoder~\cite{liu2024rtlcoder}, its supervised
LoRA adapter, and the subsequent correctness-gated policy adapter.  Each policy
is sampled \claim{SamplesPerDesign} times per primary design.  A candidate is
accepted only if it matches the reference under \claim{EvalOracleStreams}
held-out stimulus streams, each containing \claim{EvalVectorsPerStream}
vectors.  Distinct accepted RTL is stored once together with its original
multiplicity.

The second-backbone replication uses Qwen-Coder and
\claim{QwenSamplesPerDesign} samples per policy on
\claim{QwenDesignCount} designs from the earlier
\claim{QwenFamilyCount}-family split.  It is reported separately because it
does not include the later recursive and median families.

\subsection{Physical evaluation}

Every reported candidate frequency comes from synthesis, placement, and routing
for part \texttt{\claim{TargetPart}} under a requested
\claim{TargetPeriod}-ns clock.  Frequency is derived from that request and the
reported worst negative slack.  We call this ``timing-derived frequency''
because the flow does not rerun implementation at every candidate clock.  Both
policies use the same flow and constraint.

The primary endpoint is penalized equal-sample frequency as defined in
Section~\ref{sec:method}.  Secondary diagnostics are conditional frequency,
oracle correctness, family effects, candidate concentration, and exact
perfect-selector curves.  Uncertainty intervals resample paired designs within
each regime.  They are post-hoc uncertainty analyses of frozen artifacts, not
preregistered hypothesis tests.

\subsection{Contextual baselines and controls}

The stored API baseline applies the same specifications and oracle to two prompt
arms, using \claim{FrontierSamplesPerDesign} samples per arm and design.  Its
smaller budget and single API configuration make it diagnostic rather than a
general ranking of frontier models.  The HLS comparison uses post-implementation
throughput and distinguishes no-pragma C++ from engineer-placed pipeline
pragmas.  Because the RTL column selects the fastest measured optimized
candidate, HLS is a selected quality bound rather than a policy-average result.
VerilogEval supplies a separate general-RTL regression control.

\subsection{Symmetric silicon protocol}

The historical board table compared a count-weighted supervised median against
an optimized-policy maximum and is excluded from quantitative claims.  The
replacement protocol applies the same count-weighted-median candidate rule to
both policies on the pre-existing design set.  Both members of every pair share
the same golden, occupy one bitstream with a co-resident canary, and use the same
clock sweep.  A silicon number enters the manuscript only after the reserved
live-sweep artifact exists.
